% Copyright (C) 2012 Raniere Silva
% 
% This file is part of 'CNPq 126874/2012-3'.
% 
% 'CNPq 126874/2012-3' is licensed under the Creative Commons
% Attribution 3.0 Unported License. To view a copy of this license,
% visit http://creativecommons.org/licenses/by/3.0/.
% 
% 'CNPq 126874/2012-3' is distributed in the hope that it will be
% useful, but WITHOUT ANY WARRANTY; without even the implied warranty of
% MERCHANTABILITY or FITNESS FOR A PARTICULAR PURPOSE.

\documentclass[12pt,a4paper]{article}
% Filename: packages.tex
% This code is part of 'CNPq 126874/2012-3'.
% 
% Description: Pacotes utilizados.
% 
% Created: 20.08.12 07:27:58 PM
% Last Change: 20.08.12 07:27:58 PM
% 
% Author: Raniere Silva, <r.gaia.cs@gmail.com>
% 
% Copyright (c) 2012, Raniere Silva. All rights reserved.
% 
% This file is license under the terms of the GNU General Public License as published by the Free Software Foundation, either version 3 of the License, or (at your option) any later version. More details at <http://www.gnu.org/licenses/>
%
%Pacotes utilizados
\usepackage[top=3cm,left=2cm,right=2cm,bottom=3cm]{geometry} %Borda das p\'{a}gina
\usepackage[utf8]{inputenc} %Ser\'{a} utilizado a codifica\c{c}ao UTF8
\usepackage[T1]{fontenc}
%\usepackage[latin1]{inputenc} % Ser\'{a} utilizado a codifica\c{c}ao latin1
\usepackage[brazil]{babel} %Ser\'{a} utilizado o idioma portugu\^{e}s
\usepackage{indentfirst} %Identa\c{c}\~{a}o de linha
\usepackage{setspace} %Espa\c{c}amento entre linas
\usepackage{amsmath, amsfonts, amssymb, amsthm} % Pacote matem\'{a}tico
\usepackage{graphicx} %Pacote para incluir figuras
\usepackage{color} %Pacote para cores
\usepackage{tikz} %Pacote para desenho de figuras
\usepackage{pgfplots} %Pacote para desenho de gr\'{a}ficos
\usepackage{subfigure} % Pacote para subfiguras
\usepackage{url} % Pacote para url
\usepackage{hyperref} % Pacote para hyperlink
\usepackage{algorithmic} % Pacote para algoritmos
\usepackage{algorithm} % Pacote para algoritmos
\usepackage{listings} % Pacote para c\'{o}digos

%Personalização
%\pgfplotsset{
%height=0.35\textheight,
%width=0.8\textwidth,
%legend style={
%at={(0.3,0.97)},
%anchor=north,
%font=\footnotesize
%},
%xlabel={},
%ylabel={},
%xtick={},
%grid=major,
%}

\hypersetup{
%colorlinks = true,
}

\floatname{algorithm}{Algoritmo}

\algsetup{linenosize=\small}
\renewcommand{\algorithmicrequire}{\textbf{Entrada:}}
\renewcommand{\algorithmicensure}{\textbf{Saída:}}
\renewcommand{\algorithmicend}{\textbf{fim}}
\renewcommand{\algorithmicif}{\textbf{se}}
\renewcommand{\algorithmicthen}{\textbf{ent\~{a}o}}
\renewcommand{\algorithmicelse}{\textbf{caso contr\'{a}rio}}
\renewcommand{\algorithmicendif}{\algorithmicend}
\renewcommand{\algorithmicfor}{\textbf{para}}
\renewcommand{\algorithmicforall}{\textbf{para todo}}
\renewcommand{\algorithmicdo}{\textbf{fa\c{c}a}}
\renewcommand{\algorithmicendfor}{\algorithmicend}
\renewcommand{\algorithmicwhile}{\textbf{enquanto}}
\renewcommand{\algorithmicendwhile}{\algorithmicend}
\renewcommand{\algorithmicrepeat}{\textbf{repita}}
\renewcommand{\algorithmicuntil}{\textbf{at\'{e}}}
\renewcommand{\algorithmicreturn}{\textbf{retorne}}
\renewcommand{\algorithmiccomment}[1]{\hspace{2em}/* #1 */}

\renewcommand{\lstlistingname}{C\'{o}digo}
\lstset{ %
% language=Octave,                % the language of the code
basicstyle=\ttfamily\small,       % the size of the fonts that are used for the code
numbers=left,                   % where to put the line-numbers
numberstyle=\footnotesize,      % the size of the fonts that are used for the line-numbers
stepnumber=5,                   % the step between two line-numbers. If it's 1, each line 
% will be numbered
numbersep=5pt,                  % how far the line-numbers are from the code
% backgroundcolor=\color{white},  % choose the background color. You must add \usepackage{color}
showspaces=false,               % show spaces adding particular underscores
showstringspaces=false,         % underline spaces within strings
showtabs=false,                 % show tabs within strings adding particular underscores
% frame=single,                   % adds a frame around the code
tabsize=4,                      % sets default tabsize to 2 spaces
captionpos=t,                   % sets the caption-position to bottom
breaklines=true,                % sets automatic line breaking
breakatwhitespace=false,        % sets if automatic breaks should only happen at whitespace
caption={\texttt{\lstname}},                 % show the filename of files included with \lstinputlisting;
% also try caption instead of title
% escapeinside={\%*}{*)},         % if you want to add a comment within your code
% morekeywords={#}            % if you want to add more keywords to the set
}

\newtheorem{defi}{Definição}
\newtheorem{prop}{Proposição}

\begin{document}
% Identificação
%     Projeto
%     Bolsista / RA
%     Orientador
%     Local de execução
%     Vigência
\title{Implementa\c{c}\~{a}o eficiente da heur\'{i}stica de reordenamento de Cuthill-McKee reversa}
\author{Raniere Silva \\ Orientando \and Aurelio Oliveira \\ Orientador}
\maketitle

\tableofcontents

% Introdução
%     Introdução ao assunto: deve ser bastante geral
%     Informações da literatura: tornam a introdução mais específica ao assunto
%     Colocação da questão estudada: especificar os objetivos do trabalho
%     Atividades desenvolvidas: dar uma idéia geral de como foi desenvolvido o
%     trabalho
% Copyright (C) 2012 Raniere Silva
% 
% This file is part of 'CNPq 126874/2012-3'.
% 
% 'CNPq 126874/2012-3' is licensed under the Creative Commons
% Attribution 3.0 Unported License. To view a copy of this license,
% visit http://creativecommons.org/licenses/by/3.0/.
% 
% 'CNPq 126874/2012-3' is distributed in the hope that it will be
% useful, but WITHOUT ANY WARRANTY; without even the implied warranty of
% MERCHANTABILITY or FITNESS FOR A PARTICULAR PURPOSE.

\section{Método Preditor-Corretor}
Consideremos o problema de programação linear na forma padrão:
\begin{align*}
    \text{minimizar } & c^t x \\
    \text{sujeito a } & A x = b, \\
    & x \geq 0,
\end{align*}
onde $A \in \mathbb{R}^{m \times n}$ é uma matriz de posto completo $m$ e $c$,
$b$ e $x$ são vetores colunas de dimensão apropriada. Associado a este problema
temos o problema dual:
\begin{align*}
    \text{maximizar } & b^t y \\
    \text{sujeito a } & A^t y + z = c, \\
    & z \geq 0,
\end{align*}
onde $y$ é um vetor coluna de dimensão $m$ de variáveis livres e $z$ é o vetor
coluna de dimensão $n$ de variáveis de folga duais. O \textit{gap} dual é dado
por $\gamma = c^t x - b^t y$ que se reduz a $\gamma = x^t z$ para pontos primais
e duais factíveis.

A direção afim nos métodos de pontos interiores primais-duais é dada por:
\begin{align}
    \begin{bmatrix}
         0 & I & A^t \\
         Z & X & 0 \\
         A & 0 & 0
     \end{bmatrix} \begin{bmatrix}
         \tilde{x} \\
         \tilde{z} \\
         \tilde{y}
     \end{bmatrix} &= \begin{bmatrix}
         r_d \\
         r_a \\
         r_p
     \end{bmatrix},
     \label{eq:primal_dual_intp:lin_system}
\end{align}
onde $X = \diag(x)$, $Z = \diag(z)$ e os respiduos são dados por $r_p = b - A
x$, $r_d = c - A^t y - z$, $r_a = - X Z e$ e $e$ representa o vetor de uns.

Eliminando as variáveis $\tilde{z}$ de \eqref{eq:primal_dual_intp:lin_system}
obtemos o sistema aumentado:
\begin{align}
    \begin{bmatrix}
        -D & A^t \\
        A & 0
    \end{bmatrix} \begin{bmatrix}
        \tilde{x} \\
        \tilde{y}
    \end{bmatrix} &= \begin{bmatrix}
        r_1 \\
        r_2
    \end{bmatrix},
    \label{eq:primal_dual_intp:aug_lin_system}
\end{align}
onde $D = X^{-1} Z$.

A forma mais utilizada para resolver \eqref{eq:primal_dual_intp:aug_lin_system}
consiste em reduzir o sistema através de eliminção das variáveis $\tilde{x}$ a
um sistema simétrico definido positivo com a matriz de equações normais $A
D^{-1} A^t$ e aplicar então a fatoração de Cholesky.

\section{Matrizes}
\subsection{Matrizes Esparsas}
Uma matriz é dita esparsa quando a maioria de seus elementos são iguais a zero, ou seja, ela possui relativamente poucos elementos não nulos. Sistemas lineares que surgem na solução de problemas reais possuem, em sua maioria, matrizes esparsas e de ordem elevada. Nestes casos, é proibitivo armazenar toda a matriz por questão de memória do computador e o que se faz na prática é guardar somente os elementos não nulos.

Uma matriz é dita simétrica quando $A_{ij} = A_{ji}$.

\begin{figure}[!htb]
    \centering
    \begin{tikzpicture}
        \matrix (A) [matrix of math nodes,%
        left delimiter  = (,%
        right delimiter = )] at (0,0)
        {%
        X & X & X & X & X & 0 & 0 \\
        X & X & 0 & 0 & 0 & X & X \\
        X & 0 & X & 0 & 0 & X & X \\
        X & 0 & 0 & X & 0 & 0 & 0 \\
        X & 0 & 0 & 0 & X & 0 & 0 \\
        0 & X & X & 0 & 0 & X & 0 \\
        0 & X & X & 0 & 0 & 0 & X \\
        };
        \node[above, shift={(0,.5)}] at (A-1-1) {$1$};
        \node[above, shift={(0,.5)}] at (A-1-2) {$2$};
        \node[above, shift={(0,.5)}] at (A-1-3) {$3$};
        \node[above, shift={(0,.5)}] at (A-1-4) {$4$};
        \node[above, shift={(0,.5)}] at (A-1-5) {$5$};
        \node[above, shift={(0,.5)}] at (A-1-6) {$6$};
        \node[above, shift={(0,.5)}] at (A-1-7) {$7$};
        \node[right, shift={(-1.5,0)}] at (A-1-1) {$1$};
        \node[right, shift={(-1.5,0)}] at (A-2-1) {$2$};
        \node[right, shift={(-1.5,0)}] at (A-3-1) {$3$};
        \node[right, shift={(-1.5,0)}] at (A-4-1) {$4$};
        \node[right, shift={(-1.5,0)}] at (A-5-1) {$5$};
        \node[right, shift={(-1.5,0)}] at (A-6-1) {$6$};
        \node[right, shift={(-1.5,0)}] at (A-7-1) {$7$};

        \matrix (B) [matrix of math nodes,%
        left delimiter  = (,%
        right delimiter = )] at (8,0)
        {%
        X & 0 & X & X & 0 & 0 & 0 \\
        0 & X & X & X & 0 & 0 & 0 \\
        X & X & X & 0 & X & 0 & 0 \\
        X & X & 0 & X & X & 0 & 0 \\
        0 & 0 & X & X & X & X & X \\
        0 & 0 & 0 & 0 & X & X & 0 \\
        0 & 0 & 0 & 0 & X & 0 & X \\
        };
        \node[above, shift={(0,.5)}] at (B-1-1) {$1$};
        \node[above, shift={(0,.5)}] at (B-1-2) {$2$};
        \node[above, shift={(0,.5)}] at (B-1-3) {$3$};
        \node[above, shift={(0,.5)}] at (B-1-4) {$4$};
        \node[above, shift={(0,.5)}] at (B-1-5) {$5$};
        \node[above, shift={(0,.5)}] at (B-1-6) {$6$};
        \node[above, shift={(0,.5)}] at (B-1-7) {$7$};
        \node[right, shift={(-1.5,0)}] at (B-1-1) {$1$};
        \node[right, shift={(-1.5,0)}] at (B-2-1) {$2$};
        \node[right, shift={(-1.5,0)}] at (B-3-1) {$3$};
        \node[right, shift={(-1.5,0)}] at (B-4-1) {$4$};
        \node[right, shift={(-1.5,0)}] at (B-5-1) {$5$};
        \node[right, shift={(-1.5,0)}] at (B-6-1) {$6$};
        \node[right, shift={(-1.5,0)}] at (B-7-1) {$7$};
    \end{tikzpicture}
    \caption{Exemplos de matriz simétrica.}
    \label{fig:exem_matriz_simetrica}
\end{figure}

A largura de banda, $\beta$, de uma matriz simétrica $A \in \mathbb{R}^{n \times n}$ é dada pela maior distância de um elemento não nulo à diagonal principal, ou seja:
\begin{align*}
    \beta(A) &= \max_{a_{ij} \neq 0} | i - j |.
\end{align*}

\begin{figure}[!htb]
    \centering
    \begin{tikzpicture}
        \matrix (A) [matrix of math nodes,%
        left delimiter  = (,%
        right delimiter = )] at (0,0)
        {%
        X & X & X & X & X & 0 & 0 \\
        X & X & 0 & 0 & 0 & X & X \\
        X & 0 & X & 0 & 0 & X & X \\
        X & 0 & 0 & X & 0 & 0 & 0 \\
        X & 0 & 0 & 0 & X & 0 & 0 \\
        0 & X & X & 0 & 0 & X & 0 \\
        0 & X & X & 0 & 0 & 0 & X \\
        };
        \node[above, shift={(0,.5)}] at (A-1-1) {$1$};
        \node[above, shift={(0,.5)}] at (A-1-2) {$2$};
        \node[above, shift={(0,.5)}] at (A-1-3) {$3$};
        \node[above, shift={(0,.5)}] at (A-1-4) {$4$};
        \node[above, shift={(0,.5)}] at (A-1-5) {$5$};
        \node[above, shift={(0,.5)}] at (A-1-6) {$6$};
        \node[above, shift={(0,.5)}] at (A-1-7) {$7$};
        \node[right, shift={(-1.5,0)}] at (A-1-1) {$1$};
        \node[right, shift={(-1.5,0)}] at (A-2-1) {$2$};
        \node[right, shift={(-1.5,0)}] at (A-3-1) {$3$};
        \node[right, shift={(-1.5,0)}] at (A-4-1) {$4$};
        \node[right, shift={(-1.5,0)}] at (A-5-1) {$5$};
        \node[right, shift={(-1.5,0)}] at (A-6-1) {$6$};
        \node[right, shift={(-1.5,0)}] at (A-7-1) {$7$};
        % Largura de banda
        \draw (A-1-6.north east) -- (A-2-7.north east);
        \draw (A-6-1.south west) -- (A-7-2.south west);

        \matrix (B) [matrix of math nodes,%
        left delimiter  = (,%
        right delimiter = )] at (8,0)
        {%
        X & 0 & X & X & 0 & 0 & 0 \\
        0 & X & X & X & 0 & 0 & 0 \\
        X & X & X & 0 & X & 0 & 0 \\
        X & X & 0 & X & X & 0 & 0 \\
        0 & 0 & X & X & X & X & X \\
        0 & 0 & 0 & 0 & X & X & 0 \\
        0 & 0 & 0 & 0 & X & 0 & X \\
        };
        \node[above, shift={(0,.5)}] at (B-1-1) {$1$};
        \node[above, shift={(0,.5)}] at (B-1-2) {$2$};
        \node[above, shift={(0,.5)}] at (B-1-3) {$3$};
        \node[above, shift={(0,.5)}] at (B-1-4) {$4$};
        \node[above, shift={(0,.5)}] at (B-1-5) {$5$};
        \node[above, shift={(0,.5)}] at (B-1-6) {$6$};
        \node[above, shift={(0,.5)}] at (B-1-7) {$7$};
        \node[right, shift={(-1.5,0)}] at (B-1-1) {$1$};
        \node[right, shift={(-1.5,0)}] at (B-2-1) {$2$};
        \node[right, shift={(-1.5,0)}] at (B-3-1) {$3$};
        \node[right, shift={(-1.5,0)}] at (B-4-1) {$4$};
        \node[right, shift={(-1.5,0)}] at (B-5-1) {$5$};
        \node[right, shift={(-1.5,0)}] at (B-6-1) {$6$};
        \node[right, shift={(-1.5,0)}] at (B-7-1) {$7$};
        % Largura de banda
        \draw (B-1-4.north east) -- (B-4-7.north east);
        \draw (B-4-1.south west) -- (B-7-4.south west);
    \end{tikzpicture}
    \caption{Exemplos da largura de banda para matriz simétrica.}
    \label{fig:exem_bandwidth}
\end{figure}
% TODO Escrever melhor
É possível ``generalizar'' a largura de banda para cada linha da matiz simética

O envelope, $\rho$, da matriz é dado por
\begin{align*}
    \rho(A) = \sum_{i = 1}^n \beta_i(A).
\end{align*}

\begin{figure}[!htb]
    \centering
    \begin{tikzpicture}
        \matrix (A) [matrix of math nodes,%
        left delimiter  = (,%
        right delimiter = )] at (0,0)
        {%
        X & X & X & X & X & 0 & 0 \\
        X & X & 0 & 0 & 0 & X & X \\
        X & 0 & X & 0 & 0 & X & X \\
        X & 0 & 0 & X & 0 & 0 & 0 \\
        X & 0 & 0 & 0 & X & 0 & 0 \\
        0 & X & X & 0 & 0 & X & 0 \\
        0 & X & X & 0 & 0 & 0 & X \\
        };
        \node[above, shift={(0,.5)}] at (A-1-1) {$1$};
        \node[above, shift={(0,.5)}] at (A-1-2) {$2$};
        \node[above, shift={(0,.5)}] at (A-1-3) {$3$};
        \node[above, shift={(0,.5)}] at (A-1-4) {$4$};
        \node[above, shift={(0,.5)}] at (A-1-5) {$5$};
        \node[above, shift={(0,.5)}] at (A-1-6) {$6$};
        \node[above, shift={(0,.5)}] at (A-1-7) {$7$};
        \node[right, shift={(-1.5,0)}] at (A-1-1) {$1$};
        \node[right, shift={(-1.5,0)}] at (A-2-1) {$2$};
        \node[right, shift={(-1.5,0)}] at (A-3-1) {$3$};
        \node[right, shift={(-1.5,0)}] at (A-4-1) {$4$};
        \node[right, shift={(-1.5,0)}] at (A-5-1) {$5$};
        \node[right, shift={(-1.5,0)}] at (A-6-1) {$6$};
        \node[right, shift={(-1.5,0)}] at (A-7-1) {$7$};
        % Perfil
        \draw[fill=gray, fill opacity=0.2] (A-1-1.north west) -- (A-6-1.north
        west) -- (A-6-2.north west) --
        (A-7-2.south west)
        \foreach \x in {7,6,5,4,3,2,1}{
        -- (A-\x-\x.south east) -- (A-\x-\x.north east) -- (A-\x-\x.north west)
        };

        \matrix (B) [matrix of math nodes,%
        left delimiter  = (,%
        right delimiter = )] at (8,0)
        {%
        X & 0 & X & X & 0 & 0 & 0 \\
        0 & X & X & X & 0 & 0 & 0 \\
        X & X & X & 0 & X & 0 & 0 \\
        X & X & 0 & X & X & 0 & 0 \\
        0 & 0 & X & X & X & X & X \\
        0 & 0 & 0 & 0 & X & X & 0 \\
        0 & 0 & 0 & 0 & X & 0 & X \\
        };
        \node[above, shift={(0,.5)}] at (B-1-1) {$1$};
        \node[above, shift={(0,.5)}] at (B-1-2) {$2$};
        \node[above, shift={(0,.5)}] at (B-1-3) {$3$};
        \node[above, shift={(0,.5)}] at (B-1-4) {$4$};
        \node[above, shift={(0,.5)}] at (B-1-5) {$5$};
        \node[above, shift={(0,.5)}] at (B-1-6) {$6$};
        \node[above, shift={(0,.5)}] at (B-1-7) {$7$};
        \node[right, shift={(-1.5,0)}] at (B-1-1) {$1$};
        \node[right, shift={(-1.5,0)}] at (B-2-1) {$2$};
        \node[right, shift={(-1.5,0)}] at (B-3-1) {$3$};
        \node[right, shift={(-1.5,0)}] at (B-4-1) {$4$};
        \node[right, shift={(-1.5,0)}] at (B-5-1) {$5$};
        \node[right, shift={(-1.5,0)}] at (B-6-1) {$6$};
        \node[right, shift={(-1.5,0)}] at (B-7-1) {$7$};
        % Perfil
        \draw[fill=gray, fill opacity=0.2] (B-1-1.north west) -- (B-5-1.north
        west) -- (B-5-3.north west) --
        (B-5-3.south west) -- (B-6-5.north west) -- (B-7-5.south west)
        \foreach \x in {7,6,5,4,3,2,1}{
        -- (B-\x-\x.south east) -- (B-\x-\x.north east) -- (B-\x-\x.north west)
        };
    \end{tikzpicture}
    \caption{Exemplos da largura de banda para matriz simétrica.}
    \label{fig:exem_profile}
\end{figure}

\subsection{Grafos e matrizes esparsas}
A teoria dos grafos foi identificada como uma poderosa ferramenta para a computação de matrizes esparsas quando Seymour Parter [Par61] usou grafos não direcionados para modelar a eliminação Gaussiana há mais de quarenta anos atrás. Grafos podem ser usados para modelar matrizes simítricas, não simétricas, fatorações etc. Além de tornar mais fácil a compreensão e análise de algoritmos para matrizes esparsas, eles também aumentam o escopo de manipulações destas matrizes usando algoritmos e técnicas já existentes.

Um grafo é, fundamentalmente, um modo de representar uma relação binária entre objetos. Para o propósito deste trabalho, considere um grafo $G = (V, E)$ como um conjunto de vértices $V = \{v_1, v_2, \ldots \}$ (ou nós) e um conjunto de arestas $E = \{e_1, e_2, \ldots \}$. Estas arestas são representadas por pares não ordenados, por exemplo, $e_1 = (v_1 , v_2)$.

Um conhecimento básico de teoria dos grafos é importante para o entendimento do trabalho. Conceitos mais específicos serão introduzidos quando necessários. Um resumo dos principais conceitos é dado por:
\begin{description}
    \item[Grau do vértice] número de arestas incidentes no vértice.
    \item[Vértices adjacentes]  dois vértices $v_1$ e $v_2$ são adjacentes quando possue uma aresta entre eles, ou seja, $e_1 = (v_1, v_2)$.
    \item[Caminho] sequência de arestas disjuntas $\left( (x_1, x_2), (x_2, x_3), \ldots (x_{k - 1}, x_k) \right)$.
    \item[Grafo conectado (conexo)] possui caminho entre qualquer par de vértices.
    \item[Subgrafo completo (clique)] cada vértice do subgrafo possue aresta incidente em todos os outros vértices do subgrafo.
    \item[Árvore] grafo conectado sem ciclos.
\end{description}

Assim como um grafo, uma matriz tambm descreve uma relação binária entre variáveis através de seus elementos não nulos. Uma matriz simétrica $A_{n \times n}$ pode ser modelada como um grafo $G(V, E)$, onde os $n$ vértices em $V$ representam as dimensões da matriz e existe uma aresta conectando os vértices $i$ e $j$ se $a_{ij} = 0$.

\begin{figure}[!htb]
    \centering
    \begin{tikzpicture}
        \node[draw, circle] (A1) at (0,0) {1};
        \node[draw, circle] (A5) at (-2,2) {5};
        \node[draw, circle] (A4) at (-2,-2) {4};
        \node[draw, circle] (A2) at (2,2) {2};
        \node[draw, circle] (A3) at (2,-2) {3};
        \node[draw, circle] (A6) at (2,0) {7};
        \node[draw, circle] (A7) at (4,0) {6};
        \draw (A1) -- (A5) (A1) -- (A4) (A1) -- (A2);
        \draw (A1) -- (A5) (A2) -- (A7) (A3) -- (A7);
        \draw (A2) -- (A6) (A3) -- (A6) (A1) -- (A3);

        \node[draw, circle] (B1) at (8,0) {1};
        \node[draw, circle] (B5) at (6,2) {5};
        \node[draw, circle] (B4) at (6,-2) {4};
        \node[draw, circle] (B2) at (10,2) {2};
        \node[draw, circle] (B3) at (10,-2) {3};
        \node[draw, circle] (B6) at (10,0) {7};
        \node[draw, circle] (B7) at (12,0) {6};
        \draw (B1) -- (B5) (B1) -- (B4) (B1) -- (B2);
        \draw (B1) -- (B5) (B2) -- (B7) (B3) -- (B7);
        \draw (B2) -- (B6) (B3) -- (B6) (B1) -- (B3);
    \end{tikzpicture}
    \caption{Exemplos do grafo correspondente a matriz simétrica.}
    \label{fig:exem_graph}
\end{figure}
\subsection{Fatoração de matrizes esparsas}
Considerando
\begin{align*}
    A &= \begin{bmatrix}
        1 & 1 & 1 & 1 & 1 & 0 & 0 \\
        1 & 1 & 0 & 0 & 0 & 1 & 1 \\
        1 & 0 & 1 & 0 & 0 & 1 & 1 \\
        1 & 0 & 0 & 1 & 0 & 0 & 0 \\
        1 & 0 & 0 & 0 & 1 & 0 & 0 \\
        0 & 1 & 1 & 0 & 0 & 1 & 0 \\
        0 & 1 & 1 & 0 & 0 & 0 & 1
    \end{bmatrix},
\end{align*}
verifica-se que
\begin{align*}
    U &= \begin{bmatrix}
        1 &    1 &    1 &    1 &    1 &    0 &    0 \\
        0 &   -1 &    0 &   -1 &   -1 &    1 &    1 \\
        0 &    0 &   -1 &   -1 &   -1 &    1 &    1 \\
        0 &    0 &    0 &    2 &    1 &   -2 &   -2 \\
        0 &    0 &    0 &    0 &  3/2 &   -1 &   -1 \\
        0 &    0 &    0 &    0 &    0 & -2/3 &  1/3 \\
        0 &    0 &    0 &    0 &    0 &    0 & -1/2
    \end{bmatrix},
\end{align*}
onde $A = P L U$.

Permutando linhas e colunas de $A$ tal que
\begin{align*}
    A' &= \begin{bmatrix}
        1 & 0 & 1 & 1 & 0 & 0 & 0 \\
        0 & 1 & 1 & 1 & 0 & 0 & 0 \\
        1 & 1 & 1 & 0 & 1 & 0 & 0 \\
        1 & 1 & 0 & 1 & 1 & 0 & 0 \\
        0 & 0 & 1 & 1 & 1 & 1 & 1 \\
        0 & 0 & 0 & 0 & 1 & 1 & 0 \\
        0 & 0 & 0 & 0 & 1 & 0 & 1
    \end{bmatrix},
\end{align*}
verifica-se que
\begin{align*}
    U' &= \begin{bmatrix}
        1 &    0 &    1 &    1 &    0 &    0 &    0 \\
        0 &    1 &    1 &    1 &    0 &    0 &    0 \\
        0 &    0 &   -2 &   -1 &    1 &    0 &    0 \\
        0 &    0 &    0 & -3/2 &  1/2 &    0 &    0 \\
        0 &    0 &    0 &    0 &  5/3 &    1 &    1 \\
        0 &    0 &    0 &    0 &    0 & -3/5 &  2/5 \\
        0 &    0 &    0 &    0 &    0 &    0 & -1/3
    \end{bmatrix}.
\end{align*}
onde $A' = P' L' U'$.

Comparando $U$ e $U'$ verifica-se que o número de elementos nulos na parte
triangular superior da matriz $U'$ é maior na matriz $U$. E
comparando os elementos nulos na parte triangular superior de $A'$ e $U'$
verifica-se que apenas o elemento $A'_{34}$ deixou de ser nulo
enquanto que em relação a $A$ e $U$ apenas os elementos $A_{16}$,
$A_{17}$ e $A_{23}$ permaneceram nulos.
comparamos os elementos 

Pelo exemplo anterior, verifica-se a possibilidade de pré-processar a matrix
de forma a reduzir a complexidade da resolução do sistema linear. Alguns
exemplos de pré-processamentos indicados por Ghidetti
\cite{Ghidetti:2010:ComparativoReordenamento} são a minimização da
largura de banda e a redução do envelope ou \textit{profile}. ``Esses
pré-processamentos consistem em dispor os elementos não nulos da matriz o mais
próximo possível da diagonal principal.''

``No contexto da solução de sisteas lineares via métodos diretos, a
minimização da largura da bnada pode proporcionar uma redução no
preenchimento que ocorre na decomposição. (\ldots) Contudo, atualmente
sistemas lineares de grande pore são usualmente solucionados por métodos
iterativos não estacionários que por um lado não alteram a esparsidade da
matriz, mas necessitam de critérios de convergência. Em geral, um
processo para acelerar a convergência, denominado precondicionamento, se faz
necessário. Os precondicionadores baseados na decomposição LU são
amplamente utilizados por apresentar uma taxa acentuada na aceleração da
convergência. Porém, tais operações alteram a esparsidade da matriz e
portanto um processo de reordenamento da matriz é fundamental para a
eficiência de tais precondicionadores.''

\section{Cuthill-McKee reverso}
No Método Preditor-Corretor pode-se utilizar a Decomposição de Cholesky para resolver o sistema linear
% TODO Identificar o sistema linear do Método Preditor-Corretor

Para problemas de grande porte, esse sistema linear é esparso e ao utilizar a Decomposição de Cholesky pode ocorrer o preenchimento da matriz.
% TODO Incluir exemplo.

É possivel permutar linhas e colunas do sistema linear de modo a minimizar o preenchimento decorrente da Decomposição de Choleslky.
% TODO Incluir exemplo.

Cuthill e McKee \cite{Cuthill:1969:ReducingBandwidth} propuseram um algoritmo de reordenação, cujo objetivo principal é reduzir a largura de banda de uma matriz simétrica.

% TODO Mostrar equivalência com grafos
\begin{algorithm}
    \caption{Pseudo-algoritmo RCM}
    \label{alg:rcm}
    \begin{algorithmic}
        \REQUIRE Grafo $G(A)$ e um vértice $v$ de grau mínimo.
        \ENSURE $n$, novo ordenamento dos vértices de $G(A)$.
        \STATE $p = \text{ vetor de zeros}$
        \STATE $n = \text{ vetor de zeros}$
        \STATE $i = 1$
        \STATE $f = \text{ fila vazia}$
        \STATE Enfileira $v$ em $f$
        \STATE $p_v = 1$
        \WHILE{$f$ não for vazia}
            \STATE Desenfeira $f$ em $v$
            \STATE $n_v = i$
            \STATE $i = i + 1$
            \FOR{vertice $w$ adjacente a $v$}
                \IF{$p_w == 0$}
                    \STATE Enfileira $w$ em $f$
                    \STATE $p_w = 1$
                \ENDIF
            \ENDFOR
        \ENDWHILE
    \end{algorithmic}
\end{algorithm}


% Materiais e Métodos
%     Materiais: citar os equipamentos, reagentes e outros ítens utilizados,
%     informando fabricante ou fornecedor
%     Métodos: descrever os procedimentos detalhados, que possam ser
%     reproduzidos com os materiais e equipamentos descritos
\section{Implementação do Método}
O Método Cuthill-McKee Reverso foi implementado pelo autor deste relatório como
uma nova funcionalidade para o Sage (\url{http://www.sagemath.org/}), um
sistema matemático \textit{free} e \textit{open-source} licenciado sob GPL que
combina vários pacotes \textit{open-source} em uma interface comum baseada em
Python. A implementação desenvolvida foi submetida nos seguintes
\textit{patchs}:
\begin{enumerate}
    \item Ticket \#13565
        (\url{http://trac.sagemath.org/sage_trac/ticket/13565})
    \item Ticket \#13564
        (\url{http://trac.sagemath.org/sage_trac/ticket/13564})
    \item Ticket \#13581
        (\url{http://trac.sagemath.org/sage_trac/ticket/13581})
    \item Ticket \#13584
        (\url{http://trac.sagemath.org/sage_trac/ticket/13584})
    \item Ticket \#13601
        (\url{http://trac.sagemath.org/sage_trac/ticket/13601})
\end{enumerate}

Para testar a implementação, utilizou-se parte da biblioteca de testes ``Netlib LP''
(\url{http://www.netlib.org/lp/}) e na construção da matriz $A D^{-1} A^T$ que
compõe o sistema linear a ser resolvido no Método de Pontos Interiores
Primal-Dual utilizou-se $D = I$, \textit{i.e.}, $D$ como a matriz identidade.


% Resultados
%     Descrição dos resultados: deve ser clara e objetiva, resumindo os achados
%     principais que serão detalhados em tabelas e figuras.
%     Ilustrações dos resultados: tabelas e figuras são muito importantes; seu
%     número deve ser o menor possível, e elas devem ser construídas com cuidado
%     para incluir todas as informações necessárias com clareza.
%     Tabelas: devem ser numeradas sequencialmente (Tabela 1, Tabela 2, etc).
%     Seu título deve ser informativo, colocado acima e justificado à esquerda.
%     Notas de rodapé (a, b, c...) podem ser colocados diretamente abaixo da
%     mesma.
%     Figuras (fotos, esquemas, gráficos): devem ser numeradas sequencialmente
%     (Figura 1, Figura 2, etc). Seu título deve ser informativo, colocado
%     abaixo e justificado à esquerda, descrevendo o que é mostrado.
\section{Resultados Obtidos e Conclusão}
Os dados referentes aos testes computacionais encontram-se na
Tabela~\ref{tab:resul} sendo que as colunas ``Lar. Banda'', ``Envelope''
referem-se a matriz $A D^{-1} A^T$ e ``Lar. Banda R.'' e ``Envelope R.''
referem-se a matriz $P A D^{-1} A^T P^T$ em que a matriz de permutação $P$ é
obtida pelo Método Cuthill-McKee Reverso.
\begin{table}
    \centering
    \caption{Resultados experimentais do Método Cuthill-McKee Reverso para
    a ``Netlib LP''.}
    \label{tab:resul}
    \csvautotabular{bench1@relatorio_parcial.csv}
\end{table}
\begin{table}
    \centering
    \csvautotabular{bench2@relatorio_parcial.csv}
\end{table}

A partir dos dados da Tabela~\ref{bench:resul} temos que o reordenamento obtido
através do Método Cuthill-McKee Reverso resultou na
\begin{itemize}
    \item redução da largura de banda para 71 matrizes (78,02\%),
    \item manutenção da largura de banda para 7 matrizes (7.69\%),
    \item aumenta da largura de banda para 13 matrizes (14.29\%).
\end{itemize}

Dentre as matrizes que tiveram a largura de banda reduzida, a nova largura de
banda corresponde, em média, a 60\% da largura de banda inicial. Já dentre as
matrizes que tiveram a largura de banda aumentada, a nova largura de banda foi,
em média, 24\% maior que a largura de banda inicial.

Dentre as 7 matrizes que não apresentaram alteração na largura de banda,
\begin{itemize}
    \item 1 matriz teve o envelope reduzido,
    \item 5 matrizes não apresentaram alteração no envelope, e
    \item 1 matriz teve o envelope aumentado.
\end{itemize}

Pelos resultados, concluímos ser vantajoso utilizar o reordenamento obtido pelo
Método Cuthill-McKee Reverso pois este reordenamento reduz a largura de banda
para a grande maioria dos casos, aproximadamente 80\% dos testes, e essa
redução costuma ser significativa, superior a 40\%.


% Discussão / Conclusões
%     Descrição dos dados à luz da literatura
%     Descrição de possíveis fontes de erro e seu efeito sobre os dados
%     Se seus experimentos falharam, quais as sugestões para corrigir o
%     problema?
\input{conclusoes@relatorio_parcial.tex}

% Matéria encaminhada para publicação
%     Quando houver, referir resumos ou artigos científicos publicados ou
%     encaminhados para publicação
\input{publicacoes@relatorio_parcial.tex}

% Perspectivas de continuidade ou desdobramento do trabalho
%     O projeto foi concluído ou será continuado?
\section{Trabalhos Futuros}
Investigou-se a redução da largura de banda e envelope para matrizes oriunda de
problemas de programação linear. Como trabalho futuro será investigado o
impacto desta redução ao resolver tais problemas de programação linear.


% Outras atividades de interesse universitário
%     Descrever participações em congressos, cursos extra-curriculares, etc
\section{Apresentações em Congressos}
O aluno também teve o trabalho de iniciação pré-selecionado como um dos cinco
melhores para o Prêmio de Iniciação Científica do XLV Simpósio Brasileiro de
Pesquisa Operacional que será apresentado em setembro de 2013.



% Apoio
%     Citar as agências que financiaram o projeto
% Copyright (C) 2012 Raniere Silva
% 
% This file is part of 'CNPq 126874/2012-3'.
% 
% 'CNPq 126874/2012-3' is licensed under the Creative Commons
% Attribution 3.0 Unported License. To view a copy of this license,
% visit http://creativecommons.org/licenses/by/3.0/.
% 
% 'CNPq 126874/2012-3' is distributed in the hope that it will be
% useful, but WITHOUT ANY WARRANTY; without even the implied warranty of
% MERCHANTABILITY or FITNESS FOR A PARTICULAR PURPOSE.

\section{Informa\c{c}\~{o}es}
Este trabalho foi financiado pelo Conselho Nacional de Desenvolvimento Cient\'{i}fico e Tecnol\'{o}gico pelo Processo 126874/2012-3.

Este trabalho \'{e} licenciado sob a Licen\c{c}a Creative Commons Atribui\c{c}\~{a}o 3.0 N\~{a}o Adaptada License. Para ver uma c\'{o}pia desta licen\c{c}a, visite \url{http://creativecommons.org/licenses/by/3.0/}.
\begin{center}
    \includegraphics{figures/cc-by.png}
\end{center}


% TODO Para a versao final, comentar as linhas a seguir.
\appendix
% Copyright (C) 2012 Raniere Silva
% 
% This file is part of 'CNPq 126874/2012-3'.
% 
% 'CNPq 126874/2012-3' is licensed under the Creative Commons
% Attribution 3.0 Unported License. To view a copy of this license,
% visit http://creativecommons.org/licenses/by/3.0/.
% 
% 'CNPq 126874/2012-3' is distributed in the hope that it will be
% useful, but WITHOUT ANY WARRANTY; without even the implied warranty of
% MERCHANTABILITY or FITNESS FOR A PARTICULAR PURPOSE.

\section{FAQ}
\begin{enumerate}
    \item \textbf{Justifique o problema de reduzir a largura de banda é
        convertido em encontrar uma renumeração dos vértices do grafo tal que
        a diferença entre os índices seja mínima?}
        \cite{Fernanda:2005:ReordenacaoCCCG}

    \item \textbf{Por que um nós de alta excentricidade produzem bons
        resultados?} \cite{Fernanda:2005:ReordenacaoCCCG}

    \item \textbf{Por que um limitante inferior para a medida a largura de
        banda de $P A P^t$ para qualquer permutação $P$ é dado pelo menor
        inteiro maior ou igual a $D/2$, onde $D$ é o grau máximo de
        qualquer nó do grafo $G(A)$?} \cite{Cuthill:1969:ReducingBandwidth}

        Tomando o nó de grau máximo a maneira de obter a menor largura de
        banda para ele é posicionando-o na diagonal de modo que metade dos nós
        adjacentes esteja de um lado da banda e a outra metade do outro lado.

    \item \textbf{Por que iniciar com um nó de grau mínimo é uma boa escolha?}
        \cite{Cuthill:1969:ReducingBandwidth}

    \item \textbf{Por que numerar os nós em ordem crescente de grau?}
        \cite{Cuthill:1969:ReducingBandwidth}


\end{enumerate}

% Copyright (C) 2012 Raniere Silva
% 
% This file is part of 'CNPq 126874/2012-3'.
% 
% 'CNPq 126874/2012-3' is licensed under the Creative Commons
% Attribution 3.0 Unported License. To view a copy of this license,
% visit http://creativecommons.org/licenses/by/3.0/.
% 
% 'CNPq 126874/2012-3' is distributed in the hope that it will be
% useful, but WITHOUT ANY WARRANTY; without even the implied warranty of
% MERCHANTABILITY or FITNESS FOR A PARTICULAR PURPOSE.

\section{Implementações}
A seguir enumeramos algumas das implementações disponíveis.

\subsection{C}
\subsubsection{Independente}
David Fritzsche, implementou o RCM para sua tese de mestrado (``Graph
Theoretical Methods for Preconditioners'', defendida em 2004 no
Departamento de Matemática da Bergische Universität Wuppertal. 

Disponível em \url{http://math.temple.edu/~daffi/software/rcm/}.

\subsection{C++}
\subsubsection{The Boost Graph Library (BGL)}
É uma biblioteca para trabalhar com grafos.

Encontra-se disponível em
\url{http://www.boost.org/doc/libs/1_51_0/libs/graph/doc/index.html}

\subsection{Python}
\subsubsection{PyOrder}
É uma biblioteca para reordenamento de matrizes esparsas.

Encontra-se disponível em \url{https://github.com/dpo/pyorder}.

\subsubsection{NetworkX}
É uma biblioteca para trabalhar com grafos.

Encontra-se disponível em \url{http://networkx.lanl.gov/index.html} e a
implementação do RCM em
\url{http://networkx.lanl.gov/examples/algorithms/rcm.html?highlight=cuthill}


% Bibliografia
%     Diversos formatos: definir qual o mais apropriado
%     IMPORTANTE
%     - Não liste se não citar.
%     - Não cite se não listar.
\newpage
\bibliographystyle{alpha}
\bibliography{../references}
\end{document}
